\begin{tcolorbox}[title=Box 1: CAF subpopulation classification in breast cancer]

\textbf{Major CAF subpopulations:}

\begin{enumerate}
  \item \textbf{myCAFs} (myofibroblastic CAFs)
  \begin{itemize}
    \item Markers: $\alpha$-SMA (ACTA2), COL1A1, COL3A1, FN1
    \item Location: peritumoural stroma
    \item Function: ECM production, tissue stiffness, immune exclusion
    \item Proportion: 40--60\% of CAFs
  \end{itemize}

  \item \textbf{iCAFs} (inflammatory CAFs)
  \begin{itemize}
    \item Markers: IL6, IL8, CXCL1, CCL2, CCL5
    \item Location: tumour margin, immune-rich areas
    \item Function: cytokine/chemokine secretion, immune recruitment
    \item Proportion: 20--40\% of CAFs
  \end{itemize}

  \item \textbf{apCAFs} (antigen-presenting CAFs)
  \begin{itemize}
    \item Markers: HLA-DRA, CD74, CD80, CD86
    \item Location: tumour--immune interface
    \item Function: antigen presentation to CD4+ T cells
    \item Proportion: 10--20\% of CAFs
  \end{itemize}
\end{enumerate}

\textbf{Additional subpopulations identified by spatial omics:}
\begin{itemize}
  \item ECM-myCAFs: high collagen production in peritumoural niche
  \item Contractile myCAFs: high Rho/ROCK activity at invasive front
  \item Regulatory iCAFs: expressing immune checkpoint ligands
  \item Quiescent fibroblasts: in distant stroma, maintaining homeostasis
\end{itemize}

\textbf{Key point:} CAF subpopulations form a continuum rather than discrete states, with local signalling gradients determining the balance between subtypes.

\end{tcolorbox}
